\documentclass[10pt,a4paper]{amsart}
\usepackage[margin=19mm]{geometry}
\usepackage{amsmath,amssymb,mathtools}
\usepackage[scr=rsfs,cal=euler]{mathalfa}
\usepackage{xfrac}
\usepackage[unicode,hidelinks,bookmarks=false]{hyperref}
\newcommand{\ie}{i.e.\ }
\newcommand{\cf}{cf.\ }
\newcommand{\resp}{respectively\ }
\newcommand{\Subsection}{\subsection}
\providecommand{\nobreakdash}{\nobreak\hbox{-}}
\setlength{\emergencystretch}{2em}
\multlinegap0pt
\usepackage{amssymb}
\usepackage{tikz}
\newtheorem{prop}{Proposition}
\title{Grothendieck's Equality vs Voevodsky's Equality}
\author{Thomas Eckl}
\date{Source: arXiv:2604.00747v1 (CC BY 4.0)}
\begin{document}
\maketitle

\noindent\emph{Source and licence.} Grothendieck's Equality vs Voevodsky's Equality. Version arXiv:2604.00747v1 (CC BY 4.0), distributed by arXiv under the Creative Commons Attribution 4.0 International licence, http://creativecommons.org/licenses/by/4.0/, as recorded in the arXiv licence metadata for this exact version and shown on the abstract page https://arxiv.org/abs/2604.00747v1. Full licence notice: \texttt{licenses/arxiv-2604.00747-CC-BY-4.0.txt}.\par\smallskip

\noindent\textit{Editorial scope.} Counted unit: the article's prop prop (source0.tex lines 1205--1207). The article's complete original proof of it is delivered with it (source0.tex lines 1208--1231, one \texttt{proof} environment of 24 lines). No mathematics is written, repaired or re-derived: the statement and the proof are the article's own lines, in the article's own notation, and no retained line is substituted or re-flowed.  No further source-local context is needed: every cross-reference of every retained line resolves inside the two delivered spans.  Nothing was omitted from any retained span, so the package carries no omission record.  No retained line contains a citation, so the wrapper carries no bibliography and no bibliography entry.  No retained line contains a figure environment, an image-inclusion command or drawing code, so no image is delivered and the package declares no asset dependency.  Editorial formatting only, in addition: an AMS \texttt{amsart} preamble that restates the article's own declarations and macros used by the retained lines, namely the environments \texttt{prop} on separate counters, together with the standard packages those lines need.  The retained environments print in this document as Proposition 1 in order of appearance, whereas the article prints the same items with the numbering of its own full document.  The complete native source of the article as distributed in the arXiv e-print for this exact version is bundled unchanged as \texttt{sources/197609/source0.tex}.
\begin{prop}
Declaring $(r, s)$ and $(r^\prime, s^\prime)$ to be related if there (merely) is a $t \in S$ such that $t(rs^\prime - r^\prime s) = 0$ defines an equivalence relation on $R \times S$ whose set quotient $R\{S^{-1}\}$ together with the map $q : R \rightarrow R\{S^{-1}\}$ given by $r \mapsto \frac{r}{1}$ (where $\frac{r}{s}$ denotes the equivalence class of the pair $(r,s)$) is a localisation of $R$ in $S$.
\end{prop}

\begin{proof}
Reflexivity and symmetry of the relation are immediate. For the transitivity, assume that $(r,s) \sim (r^\prime,s^\prime)$ and $(r^\prime,s^\prime) \sim (r^{\prime\prime},s^{\prime\prime})$ or equivalently, $t(rs^\prime - r^\prime s) = 0$ and $t^\prime(r^\prime s^{\prime\prime} - r^{\prime\prime} s^\prime) = 0$ for some $t, t^\prime \in S$. Then
\[ tt^\prime s^\prime(r s^{\prime\prime} - r^{\prime\prime} s) = t^\prime s^{\prime\prime} \cdot t(rs^\prime - r^\prime s) + ts \cdot t^\prime(r^\prime s^{\prime\prime} - r^{\prime\prime} s^\prime) = 0, \]
so $(r,s) \sim (r^{\prime\prime},s^{\prime\prime})$, since $tt^\prime s^\prime$ is an element in the monoid $S$.

The usual calculation rules for fractions, $(r,s) + (r^\prime, s^\prime) = (r s^\prime + r^\prime s, ss^\prime)$ and $(r,s) \times (r^\prime, s^\prime) = (rr^\prime, ss^\prime)$ provide addition and multiplication on the set quotient, as these operations preserve the equivalence relation: If $(r,s) \sim (q,t)$, that is, there is $u \in S$ such that $u(rt - qs) = 0$, then 
\[ (r,s) + (r^\prime, s^\prime) = (rs^\prime + r^\prime s, ss^\prime) \sim (qs^\prime + r^\prime t, ts^\prime) = (q,t) + (r^\prime, s^\prime) \]
because 
\[ u((rs^\prime + r^\prime s)ts^\prime - (qs^\prime + r^\prime t)ss^\prime) = u(s^\prime)^2(rt - qs) = 0, \]
and 
\[ (r,s) \times (r^\prime, s^\prime) = (rr^\prime, ss^\prime) \sim (qr^\prime, ts^\prime) = (q,t) \times (r^\prime, s^\prime) \]
because
\[ u(rr^\prime ts^\prime - qr^\prime ss^\prime) = r^\prime s^\prime (rt - qs) = 0. \]
It is straightforward to verify additive and multiplicative associativity and commutativity. $0$ will be the equivalence class of $(0,1)$, and $1$ will be the equivalence class of $(1,1)$. The negative of (the equivalence class of) $(r,s)$ will be (the equivalence class of) $(-r,s)$: 
\[ (r,s) + (-r,s) = (rs - rs,s^2) = (0,s^2) \sim (0,1). \]
This is actually the only law which fails when trying to define the ring structure directly on $R \times S$ - this is only possible on the quotient $R\{S^{-1}\}$.

The map $R \rightarrow R\{S^{-1}\}$ given by $r \mapsto \frac{r}{1}$ is a ring homomorphism: $0$ and $1$ are mapped to $0$ and $1$, and the map is compatible with addition and multiplication.

So it remains to show that $R\{S^{-1}\}$ is a localisation. Let $f: R \rightarrow Q$ be a ring homomorphism to a commutative ring $Q$ with one that maps elements of $S$ to units of $Q$. The map (of sets) $\overline{f}: R \times S \rightarrow Q$ given by $(r,s) \mapsto f(r)f(s)^{-1}$ factorizes $f$ through the map $\lambda: r \mapsto (r,1)$. If $(r,s) \sim (r^\prime,s^\prime)$, that is, there is $u \in S$ such that $u(rs^\prime - r^\prime s) = 0$, we have $f(u)(f(r)f(s^\prime - f(r^\prime)f(s)) = 0$, and by multiplying with the inverses of $f(u)$, $f(s)$ and $f(s^\prime)$ we can conclude that
$\overline{f}(r,s) = \overline{f}(r^\prime,s^\prime)$. 

Therefore, the universal property of set quotients yields a unique map $g : R\{S^{-1}\} \rightarrow Q$ factorising the quotient map $q : R \times S \rightarrow R\{S^{-1}\}$, and it is straightforward to show that the composed map $g$ is a ring homomorphism such that $f = g \circ q \circ \lambda$. It is also the unique such ring homomorphism, as $g(\frac{r}{s}) = f(r)f(s)^{-1}$ is determined by $f(r)$ and $f(s)$.  
\end{proof}

\end{document}
